\documentclass[11pt,a4paper]{article}

% Idioma y codificación
\usepackage[spanish,es-tabla]{babel}
\usepackage[T1]{fontenc}
\usepackage[utf8]{inputenc}
\usepackage{csquotes}
\usepackage{multicol}
\usepackage{siunitx}
\sisetup{round-mode=places,round-precision=2,detect-all}

% Formato y utilidades
\usepackage[a4paper,margin=2.5cm]{geometry}
\usepackage{graphicx}
\graphicspath{{../anexos/}{figuras/}{imagenes/}{img/}}
\usepackage{float}
\usepackage{hyperref}
\usepackage{booktabs}
\hypersetup{
  colorlinks=true,
  linkcolor=black,
  urlcolor=blue,
  citecolor=black
}

% Bibliografía
\usepackage[style=ieee,backend=biber]{biblatex}
\addbibresource{referencias.bib}

\begin{document}

% ----------------- Portada -----------------
% Portada simple
\begin{titlepage}
  \centering
  {\Large Universidad Tecnol\'ogica del Uruguay (UTEC)\par}
  \vspace{0.5cm}
  {\large Proyecto Integrador — Brazo rob\'otico did\'actico\par}
  \vspace{2cm}
  {\LARGE Evidencias del proyecto\par}
  \vspace{1.5cm}
  \begin{tabular}{rl}
    Integrantes: & Mateo Lecuna, Hector Pereira, Priscila Rossi \\
    Docente/s:   & Marcelo Diaz, Christian Campero \\
    Fecha:       & \today
  \end{tabular}
  \vfill
\end{titlepage}

\tableofcontents
\newpage

% ----------------- Secciones principales -----------------
\section{Resumen}
Cuales fueron los avances principales que tuvimos

\section{Metodología y cronograma}

\section{Evidencias de diseño y fabricación}

\subsection{Materiales y métodos}

\subsection{Brazo articulado}
% Luego de diseñar las articulaciones correctamente...

% Figura del diseño y ensamblaje preliminar
% Figura del render del diseño final
% Figura del diseño final ensamblado

\subsection{Controlador espejo}
% Figura diseño original
% Explicacion diseño original
% Figura nuevo diseño
% Explicacion nuevo diseño




\section{Electrónica y programación}
% Explicacion codigo implementado hasta el momento
% Figura del circuito para el control espejo de los servomotores
% Explicacion del circuito implementado
% Diagramas
% Explicacion del cod

\newpage

% ----------------- Apéndices -----------------
\appendix
\section{Inspiraciones y referencias de mercado}

\begin{multicols}{2}
% \begin{figure}[H]
%   \centering
%   \includegraphics[width=0.4\textwidth]{ruta/imagen1.png}
%   \caption{Modelo referencial.}
% \end{figure}
\end{multicols}

\section{Código}

\begin{verbatim}
#include <avr/io.h>
#include <util/delay.h>

int main(void) {
    // Código base
    while (1) {
        // Bucle principal
    }
}
\end{verbatim}

\printbibliography

\end{document}
